\section{Validity and Scope}
This study is intentionally narrow. It supports live AKS claims about FinOps attribution, quality-gate execution, budget tracking through \texttt{Warning}, \texttt{Critical}, and \texttt{Exceeded}, declared-residency reporting plus an FR-only enforce-mode reroute, human approval gating with route actuation, true-positive shadow-AI detection on labeled rogue workloads, an eight-baseline routing campaign confounded by provider throttling, and small-to-medium performance profiles on one cluster. The compliance framing is deliberately limited to configured policy enforcement and declared-residency observability rather than formal legal certification, which is important when reading the work against GDPR, the EU AI Act, NIST AI RMF, OWASP guidance, and ISO/IEC 42001 \cite{gdprInfo,euAIActSummary,nistAIRMF,owaspLLMTop10,iso42001}. It does not yet support:
\begin{itemize}[leftmargin=*]
\item final routing-tradeoff claims against a full baseline family,
\item confidence-interval-backed or other inferential headline numbers,
\item broad residency-enforcement claims that ignore live catalog routability limits and quality-gate bypasses,
\item tamper-evident approval or audit-integrity claims,
\item full shadow-AI precision/recall characterization or adversarial evasion coverage, or
\item production-scale performance and scalability claims.
\end{itemize}

Additional validity threats are specific and concrete: all evidence comes from one cluster rebuilt on one day; the E2 campaign was dominated by Azure OpenAI 429 throttling on the France deployment; the E8 large profile was blocked by a single-node scheduler limit rather than by a multi-node stress test; the improved E3 campaign still covers only four application-specific datasets despite raising them to 20 prompts and five repetitions, and one gate (\texttt{legal-contract-quality}) remained verdict-unstable across repetitions; and the current shadow-AI evidence is true-positive only. Accordingly, the paper should be read as a measured AKS systems study with a deliberately narrowed claim set, not as a final comparative routing, security, or scalability paper.
